\import{set.tex}
\import{set/partition.tex}
\import{relation.tex}
\import{order/quasiorder.tex}


\section{Partial equivalences and equivalences}

Partial equivalences are also called restricted equivalences.

\begin{abbreviation}\label{partialequivalence}
    $E$ is a partial equivalence iff
    $E$ is a relation and $E$ is transitive and symmetric.
\end{abbreviation}

\begin{classification}\label{partialequivalence_is_relation}
    Let $E$ be a partial equivalence.
    Then $E$ is a relation.
\end{classification}
\begin{proof}
\end{proof}

\begin{proposition}\label{partialequivalence_is_quasireflexive}
    Let $E$ be a partial equivalence.
    Then $E$ is quasireflexive.
\end{proposition}
\begin{proof}
    Follows by \auto{quasireflexive,symmetric,transitive}.
\end{proof}

\begin{classification}\label{partialequivalence_is_quasiorder}
    Let $E$ be a partial equivalence.
    Then $E$ is a quasiorder.
\end{classification}
\begin{proof}
    $E$ is quasireflexive by \obv{partialequivalence_is_quasireflexive}.
    Follows by definition.
\end{proof}

\begin{abbreviation}\label{equivalence_on}
    $E$ is an equivalence on $A$ iff
    $E$ is a symmetric quasiorder on $A$.
\end{abbreviation}

\begin{proposition}\label{equivalence_on_implies_partialequivalence}
    Suppose $E$ is an equivalence on $A$.
    Then $E$ is a partial equivalence.
\end{proposition}
\begin{proof}
    $E$ is a quasiorder by \obv{quasiorder_on_implies_quasiorder}.
    $E$ is transitive by definition.
    $E$ is symmetric by definition.
    Follows by definition.
\end{proof}

\begin{proposition}\label{partialequivalence_is_equivalence_on_field}
    Let $E$ be a partial equivalence.
    Then $E$ is an equivalence on $\fld{E}$.
\end{proposition}
\begin{proof}
    $E$ is a binary relation on $\fld{E}$ by notion.
    $E$ is reflexive on $\fld{E}$ by
        \obv{partialequivalence_is_quasireflexive,
            quasireflexive_implies_reflexive_on_fld}.
    $E$ is transitive and symmetric by definition.
\end{proof}

\begin{proposition}\label{equivalence_on_iff_partialequivalence_and_field}
    $E$ is an equivalence on $A$ iff
    $E$ is a partial equivalence and $\fld{E}=A$.
\end{proposition}
\begin{proof}
    Show if $E$ is an equivalence on $A$ then
        $E$ is a partial equivalence and $\fld{E}=A$.
    \begin{subproof}
        Assume $E$ is an equivalence on $A$.
        $E$ is a partial equivalence by
            \obv{equivalence_on_implies_partialequivalence}.
        $E$ is a binary relation on $A$ by definition.
        Hence $\fld{E}\subseteq A$ by
            \obv{binary_relation_field_subseteq}.
        $E$ is reflexive on $A$ by definition.
        Hence $A\subseteq\fld{E}$ by
            \auto{reflexive_on,fld_intro_left,subseteq}.
        Thus $\fld{E}=A$ by \obv{subseteq_antisymmetric}.
    \end{subproof}
    Show if $E$ is a partial equivalence and $\fld{E}=A$ then
        $E$ is an equivalence on $A$.
    \begin{subproof}
        Assume $E$ is a partial equivalence and $\fld{E}=A$.
        $E$ is an equivalence on $\fld{E}$ by
            \obv{partialequivalence_is_equivalence_on_field}.
    \end{subproof}
    Follows by \auto.
\end{proof}

\begin{proposition}\label{inters_of_family_of_partial_equivalences_is_partial_equivalence}
    Let $F$ be a family of relations.
    Suppose every element of $F$ is a partial equivalence.
    Then $\inters{F}$ is a partial equivalence.
\end{proposition}
\begin{proof}
    $\inters{F}$ is a relation
        by \obv{inters_of_family_of_relations_is_relation}.
    $\inters{F}$ is symmetric by \auto{symmetric,inters_iff_forall,inters_destr}.
    $\inters{F}$ is transitive by \auto{transitive,inters_iff_forall,inters_destr}.
\end{proof}

\begin{proposition}\label{inters_of_family_of_equivalences_on_a_set_is_equivalence_on_that_set}
    Let $F$ be an inhabited family of relations.
    Suppose every element of $F$ is an equivalence on $A$.
    Then $\inters{F}$ is an equivalence on $A$.
\end{proposition}
\begin{proof}
    $\inters{F}$ is reflexive on $A$ by \auto{inters_of_family_of_reflexive_relations_is_reflexive}.
    $\inters{F}$ is a partial equivalence by
        \auto{equivalence_on_implies_partialequivalence,
            inters_of_family_of_partial_equivalences_is_partial_equivalence}.
    $\inters{F}$ is a binary relation on $A$ by
        \auto{inters_subseteq_elem,subseteq_transitive}.
\end{proof}
% NOTE: the union of a family of equivalence relations is not necessarily an
% equivalence relation because a binary union of transitive relations is not
% necessarily transitive.


\subsection{Equivalence classes}

\begin{abbreviation}\label{equiv_class}
    $\equivalenceClass{E}{a} = \downward{E}{a}$.
\end{abbreviation}

\begin{abbreviation}\label{equivclass_abbr}
    The $E$-equivalence class of $a$ is $\equivalenceClass{E}{a}$.
\end{abbreviation}

\begin{proposition}\label{partialequivalence_equivclass_inhabited}
    Let $E$ be a partial equivalence.
    Let $a\in \fld{E}$.
    Then $a\in\equivalenceClass{E}{a}$.
\end{proposition}
\begin{proof}
    $E$ is reflexive on $\fld{E}$ by
        \obv{partialequivalence_is_quasireflexive,
            quasireflexive_implies_reflexive_on_fld}.
    $a\mathrel{E} a$ by \obv{reflexive_on}.
    Follows by \obv{downward_closure_iff}.
\end{proof}

\begin{proposition}\label{equivclasses_inhabited_equivalenceon}
    Let $E$ be an equivalence on $A$.
    Let $a\in A$.
    Then $a\in\equivalenceClass{E}{a}$.
\end{proposition}
\begin{proof}
    $a\mathrel{E} a$ by \obv{reflexive_on}.
    Follows by \obv{downward_closure_iff}.
\end{proof}

\begin{proposition}\label{equiv_implies_equivanlence_classes_eq}
    Let $E$ be an equivalence on $A$.
    Let $a, b\in A$.
    Suppose $a\mathrel{E} b$.
    Then $\equivalenceClass{E}{a} = \equivalenceClass{E}{b}$.
\end{proposition}
\begin{proof}
    We have for all $c\in\equivalenceClass{E}{a}$
        we have $c\in\equivalenceClass{E}{b}$
        by \obv{downward_closure_iff,transitive}.
    We have for all $c\in\equivalenceClass{E}{b}$
        we have $c\in\equivalenceClass{E}{a}$
        by \auto{downward_closure_iff,symmetric,transitive}.
    Follows by set extensionality.
\end{proof}

\begin{proposition}\label{equivclasses_eq_implies_equiv}
    Let $E$ be an equivalence on $A$.
    Let $a, b\in A$.
    Suppose $\equivalenceClass{E}{a} = \equivalenceClass{E}{b}$.
    Then $a\mathrel{E} b$.
\end{proposition}
\begin{proof}
    $a\in\equivalenceClass{E}{a}$ by \obv{equivclasses_inhabited_equivalenceon}.
    Thus $a\in\equivalenceClass{E}{b}$.
    Follows by \obv{downward_closure_iff}.
\end{proof}

\begin{proposition}\label{equiv_iff_equivclasses_eq}
    Let $E$ be an equivalence on $A$.
    Let $a, b\in A$.
    Then $a\mathrel{E} b$ iff $\equivalenceClass{E}{a} = \equivalenceClass{E}{b}$.
\end{proposition}
\begin{proof}
    Follows by \auto{equiv_implies_equivanlence_classes_eq,equivclasses_eq_implies_equiv}.
\end{proof}

\begin{proposition}\label{equivclasses_diseq_implies_disjoint}
    Let $E$ be a partial equivalence.
    Suppose $\equivalenceClass{E}{a}\neq \equivalenceClass{E}{b}$.
    Then $\equivalenceClass{E}{a}$ is disjoint from $\equivalenceClass{E}{b}$.
\end{proposition}
\begin{proof}
    Show there exists no $c$ such that
        $c\in\equivalenceClass{E}{a}$ and
        $c\in\equivalenceClass{E}{b}$.
    \begin{subproof}
        Assume there exists a set $c$ such that
            $c\in\equivalenceClass{E}{a}$ and
            $c\in\equivalenceClass{E}{b}$.
        Take $c$ such that $c\in\equivalenceClass{E}{a}$ and
            $c\in\equivalenceClass{E}{b}$ by assumption.
        We have $c\mathrel{E}a$ and $c\mathrel{E}b$
            by \obv{downward_closure_iff}.
        Thus $a\mathrel{E}c$ and $b\mathrel{E}c$ by \auto{symmetric}.
        Thus $a\mathrel{E}b$ and $b\mathrel{E}a$ by \obv{transitive}.
        Hence $\equivalenceClass{E}{a}\subseteq\equivalenceClass{E}{b}$
            by \obv{downward_closure_iff,transitive_downward_subseteq}.
        We have $\equivalenceClass{E}{b}\subseteq\equivalenceClass{E}{a}$
            by \obv{downward_closure_iff,transitive_downward_subseteq}.
        Thus $\equivalenceClass{E}{a}=\equivalenceClass{E}{b}$
            by \obv{subseteq_antisymmetric}.
        Contradiction.
    \end{subproof}
    Follows by \obv{disjoint}.
\end{proof}

\begin{corollary}\label{equivalenceon_equivclasses_diseq_implies_disjoint}
    Let $E$ be an equivalence on $A$.
    Suppose $\equivalenceClass{E}{a}\neq \equivalenceClass{E}{b}$.
    Then $\equivalenceClass{E}{a}$ is disjoint from $\equivalenceClass{E}{b}$.
\end{corollary}
\begin{proof}
    Follows by
        \obv{equivalence_on_implies_partialequivalence,
            equivclasses_diseq_implies_disjoint}.
\end{proof}


\subsection{Quotients}

\begin{definition}\label{quotient}
    $\quotient{A}{E} = \{\equivalenceClass{E}{a}\mid a\in A\}$.
\end{definition}

\begin{proposition}\label{quotient_emptyset}
    $\quotient{\emptyset}{\emptyset} = \emptyset$.
\end{proposition}
\begin{proof}
    We have for all $C\in\quotient{\emptyset}{\emptyset}$
        we have $C\in\emptyset$ by \auto{quotient,emptyset}.
    We have for all $C\in\emptyset$
        we have $C\in\quotient{\emptyset}{\emptyset}$ by \auto{emptyset}.
    Follows by set extensionality.
\end{proof}


\begin{proposition}\label{quotient_elems_disjoint}
    Let $E$ be an equivalence on $A$.
    Suppose $B,C\in\quotient{A}{E}$ and $B\neq C$.
    Then $B$ is disjoint from $C$.
\end{proposition}
\begin{proof}
    Take $b$ such that $b\in A$ and $B = \equivalenceClass{E}{b}$ by \obv{quotient}.
    Take $c$ such that $c\in A$ and $C = \equivalenceClass{E}{c}$ by \obv{quotient}.
    Then $B$ is disjoint from $C$ by \auto{equivalenceon_equivclasses_diseq_implies_disjoint}.
\end{proof}

\begin{proposition}\label{quotient_elems_inhabited}
    Let $E$ be an equivalence on $A$.
    Suppose $C\in\quotient{A}{E}$.
    Then $C$ is inhabited.
\end{proposition}
\begin{proof}
    Take $a$ such that $a\in A$ and $C = \equivalenceClass{E}{a}$ by \obv{quotient}.
    Then $a\in \equivalenceClass{E}{a}$ by \obv{equivclasses_inhabited_equivalenceon}.
    Follows by \auto.
\end{proof}

\begin{proposition}\label{quotient_elems_type}
    Let $E$ be an equivalence on $A$.
    Suppose $a\in C\in\quotient{A}{E}$.
    Then $a\in A$.
\end{proposition}
\begin{proof}
    Take $b$ such that $b\in A$ and $C = \equivalenceClass{E}{b}$
        by \obv{quotient}.
    Then $a\mathrel{E} b$ by \obv{downward_closure_iff}.
    Thus $a\in A$ by \obv{times_tuple_elim,subseteq}.
\end{proof}


\begin{corollary}\label{quotient_notni_emptyset}
    Let $E$ be an equivalence on $A$.
    $\emptyset\notin\quotient{A}{E}$.
\end{corollary}
\begin{proof}
    Follows by \auto{quotient_elems_inhabited,emptyset}.
\end{proof}

\begin{proposition}\label{quotient_partition}
    Let $E$ be an equivalence on $A$.
    $\quotient{A}{E}$ is a partition.
\end{proposition}
\begin{proof}
    $\emptyset\notin\quotient{A}{E}$ by \obv{quotient_notni_emptyset}.
    For all $B, C\in \quotient{A}{E}$ such that $B\neq C$
        we have $B$ is disjoint from $C$ by \obv{quotient_elems_disjoint}.
    Follows by \obv{partition}.
\end{proof}

\begin{proposition}\label{quotient_partition_of}
    Let $E$ be an equivalence on $A$.
    $\quotient{A}{E}$ is a partition of $A$.
\end{proposition}
\begin{proof}
    We have for all $a\in\unions{(\quotient{A}{E})}$
        we have $a\in A$ by \auto{unions_iff,quotient_elems_type}.
    We have for all $a\in A$
        we have $a\in\unions{(\quotient{A}{E})}$
        by \auto{unions_iff,quotient,equivclasses_inhabited_equivalenceon}.
    Thus $\unions{(\quotient{A}{E})} = A$ by set extensionality.
    Follows by \obv{quotient_partition}.
\end{proof}



\begin{definition}\label{equivalence_from_partition}
    $\equivfrompartition{P}{A} = \{(a, b)\mid a\in A, b\in A\mid \exists C\in P.\ a, b\in C\}$.
\end{definition}

\begin{proposition}\label{equivalence_from_partition_intro}
    Let $P$ be a partition of $A$.
    Let $a,b\in A$.
    Suppose $a,b\in C\in P$.
    Then $a\mathrel{\equivfrompartition{P}{A}} b$.
\end{proposition}
\begin{proof}
    Follows by \auto{equivalence_from_partition}.
\end{proof}

\begin{proposition}\label{equivalence_from_partition_reflexive}
    Let $P$ be a partition of $A$.
    $\equivfrompartition{P}{A}$ is reflexive on $A$.
\end{proposition}
\begin{proof}
    Follows by \auto{reflexive_on,equivalence_from_partition_intro,partition,unions_iff}.
\end{proof}

\begin{proposition}\label{equivalence_from_partition_symmetric}
    Let $P$ be a partition.
    $\equivfrompartition{P}{A}$ is symmetric.
\end{proposition}
\begin{proof}
    It suffices to show that for all $a,b$ we have
        $a\mathrel{\equivfrompartition{P}{A}}b$ iff
        $b\mathrel{\equivfrompartition{P}{A}}a$
        by \obv{symmetric}.
    Fix $a$.
    Fix $b$.
    Show if $a\mathrel{\equivfrompartition{P}{A}}b$ then
        $b\mathrel{\equivfrompartition{P}{A}}a$.
    \begin{subproof}
        Assume $a\mathrel{\equivfrompartition{P}{A}}b$.
        $a,b\in A$ by \auto{equivalence_from_partition,pair_eq_iff}.
        Take $C\in P$ such that $a,b\in C$
            by \auto{equivalence_from_partition,pair_eq_iff}.
        Then $b,a\in C$.
        Thus $b\mathrel{\equivfrompartition{P}{A}}a$
            by \auto{equivalence_from_partition}.
    \end{subproof}
    Show if $b\mathrel{\equivfrompartition{P}{A}}a$ then
        $a\mathrel{\equivfrompartition{P}{A}}b$.
    \begin{subproof}
        Assume $b\mathrel{\equivfrompartition{P}{A}}a$.
        $b,a\in A$ by \auto{equivalence_from_partition,pair_eq_iff}.
        Take $C\in P$ such that $b,a\in C$
            by \auto{equivalence_from_partition,pair_eq_iff}.
        Then $a,b\in C$.
        Thus $a\mathrel{\equivfrompartition{P}{A}}b$
            by \auto{equivalence_from_partition}.
    \end{subproof}
    Follows by \auto.
\end{proof}

\begin{proposition}\label{equivalence_from_partition_transitive}
    Let $P$ be a partition.
    $\equivfrompartition{P}{A}$ is transitive.
\end{proposition}
\begin{proof}
    It suffices to show that for all $a,b,c$ such that
        $a\mathrel{\equivfrompartition{P}{A}}b
        \mathrel{\equivfrompartition{P}{A}}c$
        we have $a\mathrel{\equivfrompartition{P}{A}}c$
        by \obv{transitive}.
    Fix $a$.
    Fix $b$.
    Fix $c$.
    Assume $a\mathrel{\equivfrompartition{P}{A}}b
        \mathrel{\equivfrompartition{P}{A}}c$.
    $a,b\in A$ by \auto{equivalence_from_partition,pair_eq_iff}.
    $b,c\in A$ by \auto{equivalence_from_partition,pair_eq_iff}.
    Take $C\in P$ such that $a,b\in C$
        by \auto{equivalence_from_partition,pair_eq_iff}.
    Take $D\in P$ such that $b,c\in D$
        by \auto{equivalence_from_partition,pair_eq_iff}.
    Show $C=D$.
    \begin{subproof}
        Suppose not.
        $C$ is disjoint from $D$ by \obv{partition}.
        Thus $b\notin D$ by \auto{disjoint}.
        Contradiction.
    \end{subproof}
    Hence $a,c\in C$.
    Thus $a\mathrel{\equivfrompartition{P}{A}}c$
        by \auto{equivalence_from_partition}.
\end{proof}

\begin{proposition}\label{equivalence_from_partition_is_equivalence}
    Let $P$ be a partition of $A$.
    $\equivfrompartition{P}{A}$ is an equivalence on $A$.
\end{proposition}
\begin{proof}
    $\equivfrompartition{P}{A}\subseteq A\times A$
        by \auto{equivalence_from_partition,subseteq,times_tuple_intro,pair_eq_iff}.
    $\equivfrompartition{P}{A}$ is reflexive on $A$
        by \obv{equivalence_from_partition_reflexive}.
    $\equivfrompartition{P}{A}$ is transitive
        by \obv{equivalence_from_partition_transitive}.
    $\equivfrompartition{P}{A}$ is symmetric
        by \obv{equivalence_from_partition_symmetric}.
\end{proof}

\begin{proposition}\label{equivalence_from_quotient}
    Let $E$ be an equivalence on $A$.
    Then $\equivfrompartition{\quotient{A}{E}}{A} = E$.
\end{proposition}
\begin{proof}
    $E$ is a binary relation on $A$.
    Hence $E\in\rels{A}{A}$ by \obv{rels_intro}.
    Thus $E$ is a relation by \obv{rels_is_relation}.
    $\quotient{A}{E}$ is a partition of $A$
        by \obv{quotient_partition_of}.
    Hence $\equivfrompartition{\quotient{A}{E}}{A}$ is an equivalence on $A$
        by \obv{equivalence_from_partition_is_equivalence}.
    Thus $\equivfrompartition{\quotient{A}{E}}{A}$ is a binary relation on $A$.
    Hence $\equivfrompartition{\quotient{A}{E}}{A}\in\rels{A}{A}$
        by \obv{rels_intro}.
    Thus $\equivfrompartition{\quotient{A}{E}}{A}$ is a relation
        by \obv{rels_is_relation}.
    It suffices to show that for all $a,b$ we have
        $a\mathrel{\equivfrompartition{\quotient{A}{E}}{A}}b$ iff
        $a\mathrel{E}b$
        by \obv{relext}.
    Fix $a$.
    Fix $b$.
    Show if $a\mathrel{\equivfrompartition{\quotient{A}{E}}{A}}b$ then
        $a\mathrel{E}b$.
    \begin{subproof}
        Assume $a\mathrel{\equivfrompartition{\quotient{A}{E}}{A}}b$.
        Take $C\in\quotient{A}{E}$ such that $a,b\in C$
            by \auto{equivalence_from_partition,pair_eq_iff}.
        Take $c\in A$ such that $C=\equivalenceClass{E}{c}$
            by \obv{quotient}.
        Then $a\mathrel{E}c$ and $b\mathrel{E}c$
            by \auto{downward_closure_iff}.
        $E$ is symmetric and transitive.
        Thus $c\mathrel{E}b$ by \obv{symmetric}.
        Hence $a\mathrel{E}b$ by \obv{transitive}.
    \end{subproof}
    Show if $a\mathrel{E}b$ then
        $a\mathrel{\equivfrompartition{\quotient{A}{E}}{A}}b$.
    \begin{subproof}
        Assume $a\mathrel{E}b$.
        Then $(a,b)\in A\times A$ by \obv{subseteq}.
        Hence $a,b\in A$ by \obv{times_tuple_elim}.
        Let $C=\equivalenceClass{E}{b}$.
        $C\in\quotient{A}{E}$ by \obv{quotient}.
        $a\in C$ by \obv{downward_closure_iff}.
        $b\in C$ by \obv{equivclasses_inhabited_equivalenceon}.
        Thus $a\mathrel{\equivfrompartition{\quotient{A}{E}}{A}}b$
            by \obv{equivalence_from_partition_intro}.
    \end{subproof}
    Follows by \auto.
\end{proof}

\begin{lemma}\label{partition_blocks_meeting_eq}
    Let $P$ be a partition.
    Suppose $C,D\in P$.
    Suppose $c\in C,D$.
    Then $C=D$.
\end{lemma}
\begin{proof}
    Suppose not.
    $C$ is disjoint from $D$ by \obv{partition}.
    Thus $c\notin D$ by \auto{disjoint}.
    Contradiction.
\end{proof}

\begin{lemma}\label{partition_block_is_equivalence_class}
    Let $P$ be a partition of $A$.
    Suppose $C\in P$.
    Suppose $b\in C$.
    Then $\equivalenceClass{\equivfrompartition{P}{A}}{b}=C$.
\end{lemma}
\begin{proof}
    $b\in\unions{P}$ by \obv{unions_intro}.
    Hence $b\in A$ by \obv{partition}.
    It suffices to show that for all $x$ we have
        $x\in\equivalenceClass{\equivfrompartition{P}{A}}{b}$ iff $x\in C$
        by \obv{setext}.
    Fix $x$.
    Show if $x\in\equivalenceClass{\equivfrompartition{P}{A}}{b}$ then $x\in C$.
    \begin{subproof}
        Assume $x\in\equivalenceClass{\equivfrompartition{P}{A}}{b}$.
        Then $x\mathrel{\equivfrompartition{P}{A}}b$
            by \obv{downward_closure_iff}.
        Take $D\in P$ such that $x,b\in D$
            by \auto{equivalence_from_partition,pair_eq_iff}.
        $D=C$ by \obv{partition_blocks_meeting_eq}.
        Thus $x\in C$.
    \end{subproof}
    Show if $x\in C$ then
        $x\in\equivalenceClass{\equivfrompartition{P}{A}}{b}$.
    \begin{subproof}
        Assume $x\in C$.
        $x\in\unions{P}$ by \obv{unions_intro}.
        Hence $x\in A$ by \obv{partition}.
        Thus $x\mathrel{\equivfrompartition{P}{A}}b$
            by \obv{equivalence_from_partition_intro}.
        Hence $x\in\equivalenceClass{\equivfrompartition{P}{A}}{b}$
            by \obv{downward_closure_iff}.
    \end{subproof}
    Follows by \auto.
\end{proof}

\begin{proposition}\label{partition_eq_quotient_by_equivalence_from_partition}
    Let $P$ be a partition of $A$.
    Then $\quotient{A}{\equivfrompartition{P}{A}} = P$.
\end{proposition}
\begin{proof}
    It suffices to show that for all $C$ we have
        $C\in\quotient{A}{\equivfrompartition{P}{A}}$ iff $C\in P$
        by \obv{setext}.
    Fix $C$.
    Show if $C\in\quotient{A}{\equivfrompartition{P}{A}}$ then $C\in P$.
    \begin{subproof}
        Assume $C\in\quotient{A}{\equivfrompartition{P}{A}}$.
        Take $b\in A$ such that
            $C=\equivalenceClass{\equivfrompartition{P}{A}}{b}$
            by \obv{quotient}.
        Then $b\in\unions{P}$ by \obv{partition}.
        Take $D\in P$ such that $b\in D$ by \obv{unions_iff}.
        $\equivalenceClass{\equivfrompartition{P}{A}}{b}=D$
            by \obv{partition_block_is_equivalence_class}.
        Thus $C=D$.
        Hence $C\in P$.
    \end{subproof}
    Show if $C\in P$ then
        $C\in\quotient{A}{\equivfrompartition{P}{A}}$.
    \begin{subproof}
        Assume $C\in P$.
        $C\neq\emptyset$ by \obv{partition}.
        Take $b$ such that $b\in C$ by \obv{nonempty_inhabited}.
        $b\in\unions{P}$ by \obv{unions_intro}.
        Hence $b\in A$ by \obv{partition}.
        Let $D=\equivalenceClass{\equivfrompartition{P}{A}}{b}$.
        $D\in\quotient{A}{\equivfrompartition{P}{A}}$ by \obv{quotient}.
        $D=C$ by \obv{partition_block_is_equivalence_class}.
        Thus $C\in\quotient{A}{\equivfrompartition{P}{A}}$.
    \end{subproof}
    Follows by \auto.
\end{proof}
